% =====================================================================
\chapter{Combinaciones, números combinatorios y doble conteo}
\label{tema:2}
% =====================================================================

Este tema cuenta \emph{selecciones no ordenadas}: subconjuntos, multiconjuntos y repartos. Aparecen los números combinatorios, el binomio y el multinomio de Newton, la técnica del doble conteo y la forma general del principio de inclusión--exclusión. Se siguen las secciones §14.5 a §14.7 de Lehman, Leighton y Meyer~\cite{mcs}; las secciones §1.4 a §1.6 de Ross~\cite{ross}; las secciones §1.4 y §1.5, «Story proofs», de Blitzstein y Hwang~\cite{bh}, cuyo estilo de «demostración mediante una historia» se adopta aquí, y la sección~3.2 de Grinstead y Snell~\cite{gs}.

\section{Combinaciones y números combinatorios}

\begin{definicion}[Combinaciones]
Sea $A$ un conjunto con $n$ elementos y $0\le k\le n$. Una \emph{combinación} (sin repetición) de orden $k$ de $A$ es un subconjunto de $A$ con $k$ elementos. Su número se denota $\binom{n}{k}$ y se llama \emph{número combinatorio} («$n$ sobre $k$»). Por convenio, $\binom{n}{k}=0$ si $k<0$ o $k>n$.
\end{definicion}

\begin{teorema}\label{teo:t2-binom}
Para $0\le k\le n$,
\[
\binom{n}{k}=\frac{n^{\underline{k}}}{k!}=\frac{n!}{k!\,(n-k)!}.
\]
\end{teorema}

\begin{proposicion}[Propiedades básicas]\label{prop:t2-props}
Para enteros $n\ge1$ y $k$:
\begin{enumerate}
  \item simetría: $\binom{n}{k}=\binom{n}{n-k}$;
  \item fórmula de Pascal: $\binom{n}{k}=\binom{n-1}{k}+\binom{n-1}{k-1}$;
  \item absorción: $k\binom{n}{k}=n\binom{n-1}{k-1}$.
\end{enumerate}
\end{proposicion}

La fórmula de Pascal permite construir el \emph{triángulo de Pascal}, en el que cada número es la suma de los dos que tiene encima:
\[
\begin{array}{c}
1\\
1\quad 1\\
1\quad 2\quad 1\\
1\quad 3\quad 3\quad 1\\
1\quad 4\quad 6\quad 4\quad 1\\
1\quad 5\quad 10\quad 10\quad 5\quad 1
\end{array}
\]

\section{Binomio y multinomio de Newton}

\begin{teorema}[Binomio de Newton]\label{teo:t2-binomio}
Para $x,y$ reales y $n\ge0$ entero,
\[
(x+y)^{n}=\sum_{k=0}^{n}\binom{n}{k}x^{k}y^{n-k}.
\]
\end{teorema}

El módulo~0.3 demuestra este teorema por inducción; la demostración anterior es la combinatoria, que explica \emph{por qué} aparecen los números combinatorios.

\begin{corolario}\label{cor:t2-sumas}
Para $n\ge 0$, $\sum_{k=0}^{n}\binom{n}{k}=2^{n}$, y para $n\ge1$, $\sum_{k=0}^{n}(-1)^{k}\binom{n}{k}=0$.
\end{corolario}

\begin{definicion}[Coeficiente multinomial]
Sean $n=k_1+\dots+k_r$ con $k_i\ge0$ enteros. Se define
\[
\binom{n}{k_1,k_2,\dots,k_r}=\frac{n!}{k_1!\,k_2!\cdots k_r!}.
\]
\end{definicion}

\begin{teorema}[Permutaciones con repetición]\label{teo:t2-multinomial}
El número de sucesiones de longitud $n$ formadas con $k_1$ copias del símbolo $s_1$, $k_2$ del símbolo $s_2$, \dots, $k_r$ del símbolo $s_r$ es $\binom{n}{k_1,\dots,k_r}$. Equivalentemente, es el número de formas de repartir $n$ objetos distintos en $r$ cajas distinguibles, con $k_i$ objetos en la caja $i$.
\end{teorema}

\addtocounter{definicion}{1}%

\begin{teorema}[Multinomio de Newton]\label{teo:t2-multinomio}
Para reales $x_1,\dots,x_r$ y $n\ge0$,
\[
(x_1+\dots+x_r)^{n}=\sum_{\substack{k_1+\dots+k_r=n\\ k_i\ge0}}\binom{n}{k_1,\dots,k_r}x_1^{k_1}\cdots x_r^{k_r}.
\]
\end{teorema}

\section{Combinaciones con repetición}

\begin{definicion}
Una \emph{combinación con repetición} de orden $k$ de un conjunto de $n$ elementos es una selección no ordenada de $k$ elementos en la que se permite repetir; equivale a un vector $(x_1,\dots,x_n)$ de enteros $x_i\ge 0$ con $x_1+\dots+x_n=k$, donde $x_i$ es el número de veces que se elige el elemento $i$.
\end{definicion}

\begin{teorema}[Estrellas y barras]\label{teo:t2-estrellas}
El número de soluciones enteras de $x_1+\dots+x_n=k$ con todos los $x_i\ge0$ es $\binom{n+k-1}{k}=\binom{n+k-1}{n-1}$. El número de soluciones con todos los $x_i\ge1$ es $\binom{k-1}{n-1}$ (para $k\ge n\ge 1$).
\end{teorema}

\begin{table}[h]
\centering
\begin{tabular}{@{}lcc@{}}
\toprule
Selección de $k$ entre $n$ & sin repetición & con repetición \\
\midrule
ordenada (importa el orden) & $n^{\underline{k}}=\dfrac{n!}{(n-k)!}$ & $n^{k}$ \\[8pt]
no ordenada & $\dbinom{n}{k}$ & $\dbinom{n+k-1}{k}$ \\
\bottomrule
\end{tabular}
\caption{Los cuatro modelos básicos de selección. El caso con repetición y sin orden es el único en que la fórmula no es «evidente».}
\label{tab:t2-modelos}
\end{table}

\section{Doble conteo y demostraciones biyectivas}

Una \emph{demostración por doble conteo} establece una identidad $a=b$ mostrando que ambos lados cuentan el mismo conjunto de dos maneras distintas; una \emph{demostración biyectiva} la establece construyendo una biyección entre un conjunto de $a$ elementos y otro de $b$ elementos. Ya se han usado ambas en la \cref{prop:t2-props}. El esquema conviene escribirlo siempre de forma explícita: (1) describir el conjunto que se cuenta (una «historia»); (2) contarlo de una forma, obteniendo el lado izquierdo; (3) contarlo de otra, obteniendo el lado derecho~\cite[§1.5]{bh}.

\addtocounter{definicion}{1}%

\begin{proposicion}[Identidad de Vandermonde]\label{prop:t2-vandermonde}
Para enteros $m,n,k\ge0$, $\displaystyle\binom{m+n}{k}=\sum_{j=0}^{k}\binom{m}{j}\binom{n}{k-j}$.
\end{proposicion}

La demostración se pide en el \cref{eje:2-vandermonde}: un grupo de $m$ mujeres y $n$ hombres, y comités de $k$ personas clasificados por el número de mujeres.

\section{Inclusión--exclusión general}

\begin{teorema}[Inclusión--exclusión]\label{teo:t2-ie}
Para conjuntos finitos $A_1,\dots,A_m$,
\[
\Bigl|\bigcup_{i=1}^{m}A_i\Bigr|=\sum_{j=1}^{m}(-1)^{j+1}\sum_{1\le i_1<\dots<i_j\le m}\abs{A_{i_1}\cap\dots\cap A_{i_j}}.
\]
\end{teorema}

\addtocounter{definicion}{1}%

\addtocounter{definicion}{1}%

\addtocounter{definicion}{1}%

